\section{Conclusion}
\label{sec:conclusion}

\begin{table}
  \centering
  {
    \renewcommand\arraystretch{1.4}
    \setlength\tabcolsep{0.4em}
    \footnotesize
    \begin{tabular}{ccc|ccc|cc}
      \toprule
      \multicolumn{3}{c|}{Description} & \multicolumn{3}{c|}{Spaces}  & \multirow{2}{*}{NL op.} &  \multirow{2}{*}{Best attack} \\
      Name & Use Case & Ref. & Input & Key & Output \\
      \midrule
      % Alternating Moduli
      \multirow{2}{*}{Alternating} & \multirow{3}{*}{MPC}
      & \cite{TCC:BIPSW18} & $\Z_2^{384}$\dag & $\Z_2^{384}$\dag & $\Z_2$ & $\mod$ & ?? \\
      \cline{3-8}
      \multirow{2}{*}{Moduli} & &  \multirow{2}{*}{\cite{C:DGHIKS21}} & $\Z_2^{256}$\ddag & $(\Z_2^{256})^{256}$\ddag & \multirow{2}{*}{$\Z_2^{81}$} & \multirow{2}{*}{$\mod$} & ?? \\
      & &  & $\Z_2^{320}$ & $(\Z_2^{320})^{320}$ & & & ?? \\
      \midrule
      % Goldreich
      \multirow{1}{*}{Goldreich's} & \multirow{2}{*}{??} & \multirow{3}{*}{\cite{DBLP:journals/tit/YangGJL22}} & $(A_5^{1024})^{1176}$ & $\{0,1\}^{1024}$ & $\{0,1\}^{1176}$ & \multirow{3}{*}{$P_5$} & ?? \\ % m=n^1.02
      \multirow{1}{*}{PRG} & & 
      & $(A_5^{2048})^{4389}$ & $\{0,1\}^{2048}$ & $\{0,1\}^{4389}$ & & ?? \\ % m=n^1.10
      & & & $(A_5^{4096})^{19893}$ & $\{0,1\}^{4096}$ & $\{0,1\}^{19893}$ & & ?? \\ % m=n^1.19
      \cline{3-8}
      & & \cite{EC:ABBS25} & $(A_5^{8192})^{45387}$ & $\{0,1\}^{8192}$ & $\{0,1\}^{45387}$ & $P_5$ & ?? \\ % m=n^1.19
      \bottomrule
    \end{tabular}
  }
  \caption{A summary of the wPRFs we present. The ``\dag'' symbol means the primitive is now considered unsafe, while ``\ddag'' indicates an ``aggressive'' set of parameters. The data limit for all such functions is at least $2^{40}$, and the intended security level at least 128 bits. ``NL op.'' refers to the non-linear operation.}
  \label{tab:summary}
\end{table}


The goal of this Systematization of Knowledge (SoK) article was to unify and present four distinct families of weak pseudo-random functions (wPRFs) under a common lens. These families of wPRFs have emerged from the theoretical cryptography community in response to various application needs. Despite their differing origins, all the constructions we examined share a common characteristic: they are shallow, typically consisting of a single round of computation. Moreover, they operate within weakened security models when compared to standard PRFs or PRPs, notably by restricting the attacker from choosing inputs. More importantly, the inputs to these primitives  are sampled uniformly at random.

This uniform randomness of public inputs, combined with the small domain sizes typically proposed, makes differential cryptanalysis largely impractical in this context. As a result, the main attack strategies that remain applicable seem to be linear cryptanalysis, based on exploiting statistical biases in (linear combinations of) the output, and algebraic or combinatorial techniques. However, the field currently lacks a comprehensive understanding of how effective such attacks can be in practice.

In particular, a more comprehensive combinatorial analysis is needed to characterize the conditions under which an attack may succeed based on properties of the public inputs. Indeed, attacks may become feasible when public inputs fall within certain ``bad'' subsets. Identifying and understanding the structure of these subsets is an interesting and important research direction.

More broadly, there is a remarkable lack of systematic cryptanalysis in this weak-function setting. Security claims for these constructions often rely on heuristic or partial evidence, and the distinction between exponential and polynomial-time attacks becomes blurred once concrete parameters are fixed. Despite this, many of these primitives are being integrated into advanced cryptographic protocols and proposed for real-world deployment. We therefore encourage the symmetric cryptography community to analyze these constructions, and aim to facilitate this task by providing reference implementations for each of them. Targeted and rigorous cryptanalysis is essential before these primitives can be confidently used inside high-level protocols. Bridging the gap between theoretical proposals and concrete security evaluation is not just desirable,  it is necessary to ensure the security of the advanced protocols that build upon them.

%\begin{itemize}
%\item On a besoin de quantifier précisément la complexité des algorithmes d'attaque histoire de pouvoir dimensionner correctement et précisément les instances sensées offrir un niveau de sécurité donné
%\item en l'état, il n'y a pas de cryptanalyse dans le modèle ``weak nonce'' ou ``weak IV'' (demander à Anne ce qui existe), qu'on introduit
%\end{itemize}


%All constructions that we presented consist of a weakened security model than the standard one (PRF or PRP), as an attacker is not able to choose the public inputs. It is however somehow different than a known plaintext scenario, as the inputs here are taken uniformly at random.

%The fact that public inputs are taken uniformly at random, combined with the proposed size makes the differential cryptanalysis unpractical. What remains seem to be, up to now likewise linear cryptanalysis, by observing biases in (linear combination of) the output or algebraic-like attacks. \yr{C'est les seules types d'attaques partout non ?}

%Some combinatorial analysis is also missing to identify a success probability of an attack that is computable over the public inputs. Indeed, as Kerckhoff's principle states that the security should rely only on the secret of the key, this principle is not met in those constructions as it could appear than an attack could work when the public inputs lie in a specific subsets. The identification of ``bad'' public input set is thus a very interesting question that need to be answered.
%\yr{Faut-il réinsister plus ?}

%Eventually, we believe that this weak scenario drastically lacks of security analysis. Indeed, actual security relies on cryptanalysis. We believe this to not be fulfilled, specifically because when fixing parameters, the distinction between exponential/polynomial types of attacks do not make anymore sense. Nowadays, many of those proposals have a purpose to be used in practice in many advanced cryptographic protocols. Therefore a big gap is present and we think that the cryptographic community should collectively target those constructions, with careful analysis in order to gain confidence in advanced protocols before proposing them as being used and implemented.
